% Shared abstract (included by main.tex, ieee.tex, springer.tex)
Teams increasingly chain an AI model's typed decisions into a workflow,
hoping the model's reported confidence will flag the cases a human should
check. We test that hope with Jev, TypeSafe's typed decision model and
its three primitives (a categorical \choice{}, a binary \noul{} and an
ordinal \score{}), on a bug-triage workflow, and we separate two
questions that are usually merged: can a signal say which chained
decisions are wrong, and can it say which of those are worth a human's
time? We study 442 Mozilla and Linux kernel bugs, then run pre-registered
tests on 500 new bugs, four further models and an independent Kubernetes
workflow. Failure is disagreement with a target action we construct from
reference labels. A frozen decision-feature estimator ranks failures
better than final-step confidence (AUROC 0.76 vs.\ 0.55 on the new bugs,
supported for all five models), and its signal is almost entirely the
historical error rate of the chosen action. The Kubernetes replication,
on a chronological split, confirmed this (AUROC 0.78 vs.\ 0.63), as did
three of four further models, including the only other one competent at
every step. These error rates do not transfer across label mixes: fitted
on a quota-sampled development set, they can rank natural-mix errors below
chance. Re-weighting to the deployment mix repairs the ranking, but the
label-free estimate of that mix is biased and made ranking worse on
Kubernetes. Prediction did not become review value: under the
pre-registered costs no 10\% review policy beat no review, an oracle
locates the value in wrongly accepted cases the tested signals do not
isolate, and whether review pays depends on the cost ratio. Knowing which
decisions are likely wrong is not the same as knowing which are worth
correcting.
